\section{Lean formalization: exact definitions and a route to \texorpdfstring{$\Prf\,\PRAAx$}{Prf PRAAx}}

\subsection{The formal endpoint}

The repository distinguishes three formal systems that must not be conflated.  Lean's ambient type
 theory is used to define syntax and metatheorems; the ordinary first-order presentation of iPRA has
primitive-recursive terms \code{PTm}, axioms \code{PRAAx}, and derivations \code{Prf}; and the
smaller native scalar core has its own derivation relation.  The checked bridge is of the form
\begin{lstlisting}[language={}] 
theorem pra_iff_native
    {Gamma : List (Fm PTm n)} {phi : Fm PTm n} :
  (Gamma |- [PRAAx] phi) <->
    NPrf (Gamma.map NFm.ofPRA) (NFm.ofPRA phi)
\end{lstlisting}
Consequently a host theorem such as
\begin{lstlisting}[language={}] 
theorem power_cover_semantic : ...
\end{lstlisting}
is not yet an iPRA formalization.  The relevant endpoint is a term such as
\begin{lstlisting}[language={}] 
def powerCoverObjectProof :
    Prf PRAAx [] powerCoverFormula := ...
\end{lstlisting}
and only afterwards may \code{pra_iff_native} transport it to the native calculus.

The source files in this archive are blueprints.  No Lean executable or complete repository
dependency closure was available in the preparation environment, so they are not claimed to
compile.  They contain no \code{axiom}, \code{sorry}, or \code{admit}; unfinished results are named
\emph{targets}, not asserted theorems.

\subsection{Two formalization layers}

The proposed development has two deliberately separate layers.

\paragraph{Host certificate layer.}
A \code{CertSetoid} contains numerical codes, a Boolean validity predicate, a proof-relevant
relation \code{eqW x y w}, and explicit witness transformers for reflexivity, symmetry, and
transitivity.  A certified map carries equality witnesses to equality witnesses.  Certified rings
and modules add congruence transformers for their operations.  This layer is useful for testing the
interfaces and for proving generic metatheorems, but its quantifiers live in Lean rather than in the
object theory.

\paragraph{Object-iPRA layer.}
Every numerical operation is replaced by one fixed \code{PR} code.  Its application is a
\code{PTm}; every checker and invariant is an \code{Fm PTm n}; and all correctness theorems are
values of \code{Prf PRAAx}.  A theorem schema may fix the definitions of a presented ring or module
externally and remain uniform in ordinary numerical parameters.  It must not evaluate arbitrary
codes for primitive-recursive programs, since there is no primitive-recursive universal evaluator
for all such codes.

\subsection{Codecs and totalization}

All objects of the internal theory are natural numbers.  The first module should therefore supply
fixed primitive-recursive codecs for:
\begin{itemize}
\item signed integers and finite tuples;
\item lists with length, bounded lookup, update, append, map, fold, maximum, and least bounded search;
\item exponent vectors, monomials, sparse polynomials, and Laurent polynomials;
\item finite matrices and bit-mask subsets of $\{0,\ldots,n\}$;
\item ideal-combination witnesses, localization witnesses, cover certificates, and descent data.
\end{itemize}
Every operation is total.  Malformed input receives a default output together with a failure flag.
Theorems assume the finite validity predicate, so no partial function is hidden in the object
language.  The basic reconstruction lemmas have the shape
\[
  \Valid(x)=1\;\longrightarrow\;
  \operatorname{decode}(\operatorname{encode}(x))=x,
\]
and bounded access lemmas are proved by quantifier-free induction on the list prefix.

\subsection{Certified-setoid library}

For each presentation $A$, introduce fixed primitive-recursive predicates
\[
        \Valid_A(x),\qquad \EqW_A(x,y,w).
\]
The object library should construct the following derivations, uniformly in the displayed codes:
\begin{align*}
\Valid_A(x)&\to \EqW_A(x,x,\mathsf{reflW}(x)),\\
\EqW_A(x,y,w)&\to \EqW_A(y,x,\mathsf{symmW}(x,y,w)),\\
\EqW_A(x,y,u)\wedge\EqW_A(y,z,v)
 &\to \EqW_A(x,z,\mathsf{transW}(x,y,z,u,v)).
\end{align*}
For every operation $F$, add a witness transformer
\[
   \EqW_A(x,x',u)\wedge\EqW_B(y,y',v)
   \to \EqW_C(F(x,y),F(x',y'),\mathsf{mapW}(u,v)).
\]
The library never asserts $\exists w\,\EqW_A(x,y,w)$ merely because two semantic quotient classes
are equal in Lean.  The witness is part of the finite mathematical data.

A useful generic lemma is bounded-fold congruence.  Suppose two lists have the same length and each
pair of corresponding entries comes with an equality witness.  Quantifier-free induction on the
prefix produces a witness between the two additive or multiplicative folds.  This lemma will be
used repeatedly in ideal identities, multinomial expansions, and \v{C}ech differentials.

\subsection{Presented quotient rings}

For
\[
 A=\Z[X_1,\ldots,X_s]/(g_0,\ldots,g_{r-1}),
\]
a code for an element is a polynomial code.  A witness that $p=q$ is a tuple
$(h_0,\ldots,h_{r-1})$ satisfying the coefficientwise identity
\[
                     p-q=\sum_{i<r}h_i g_i.
\]
The checker expands every product and compares a bounded coefficient range determined from the
input supports.  The soundness proof uses only finite polynomial arithmetic.

The witness transformers are explicit.  Negation negates every multiplier; addition concatenates
or adds multiplier lists; and multiplication uses
\[
 pq-p'q'=(p-p')q+p'(q-q').
\]
Associativity, commutativity, distributivity, and the unit laws can be witnessed by the zero
multiplier list once the corresponding polynomial identity is normalized coefficientwise.  No
Gr\"obner basis is needed for soundness.  A Gr\"obner or elimination algorithm is a separate
\emph{producer} that may find witnesses in useful subclasses.

The principal object theorem should be exposed as a derivation of a formula asserting:
\begin{quote}
if the ideal-certificate checker accepts, then the represented quotient elements are equal, and
all quotient operations preserve accepted equality certificates.
\end{quote}

\subsection{Localization witnesses}

Fix a presented ring $A$ and an element $f$.  A fraction is a pair $(a,n)$ representing $a/f^n$.
Equality of $(a,n)$ and $(b,m)$ is witnessed by a number $k$ and a base-ring witness for
\[
                       f^k(f^m a-f^n b)=0.                 \tag{14.1}\label{eq:loc-witness-object}
\]
The object formalization should implement the following witness algebra.

\begin{itemize}
\item Reflexivity uses $k=0$ and the zero witness.
\item Symmetry negates the witness in \eqref{eq:loc-witness-object}.
\item For transitivity, suppose $k$ kills
      $f^m a-f^n b$ and $\ell$ kills $f^r b-f^m c$.  Multiplying the first
      relation by $f^{\ell+r}$ and the second by $f^{k+n}$ and adding cancels
      the middle term.  The result supplies a safe killing exponent
      $k+\ell+m$ for $f^r a-f^n c$; a slightly smaller exponent can be used
      after a more exact bookkeeping proof.
\item Addition and multiplication expand the cross differences and combine the supplied witnesses
      by finite ring identities.
\end{itemize}

Every exponentiation is a fixed primitive recursion.  The loop invariant records the current power,
the accumulated product, and an accepted equality witness.  The induction formula is
quantifier-free.

\subsection{The power-cover constructor}

The most important generic producer starts with a checked identity
\[
                         \sum_{i<r}c_i f_i=1               \tag{14.2}\label{eq:unit-cover-object}
\]
and a number $N\geq1$.  Set $M=r(N-1)+1$.  Enumerate every weak composition
$\alpha\in\N^r$ of $M$ and expand $(\sum c_i f_i)^M$.  The pigeonhole bound gives an index
$i$ with $\alpha_i\geq N$; choose the least such index and assign that monomial to the accumulator
$d_i$ after factoring $f_i^N$.  The output satisfies
\[
                         \sum_{i<r}d_i f_i^N=1.             \tag{14.3}\label{eq:power-cover-object}
\]

A direct object implementation uses nested bounded loops rather than an abstract multinomial
polynomial.  Its state contains:
\begin{enumerate}[label=(\roman*)]
\item the current exponent vector and its remaining sum;
\item the partial multinomial coefficient;
\item the chosen least large coordinate;
\item the coefficient accumulators $d_0,\ldots,d_{r-1}$;
\item the partial expansion value and the partial factored value.
\end{enumerate}
The invariant says that the two partial values are equal and that every processed composition was
assigned to exactly one accumulator.  All quantifiers are bounded folds represented by Boolean
prefix predicates.  A final finite-power calculation rewrites the expansion by \eqref{eq:unit-cover-object} to
$1^M=1$.

The endpoint should be
\begin{lstlisting}[language={}] 
def powerCoverObjectProof :
  Prf PRAAx [] powerCoverFormula := ...
\end{lstlisting}
where \code{powerCoverFormula} states validity of the input certificate, $N\geq1$, and acceptance
of the output identity checker.

\subsection{Principal-cover uniqueness and gluing}

For joint monicity, the input contains for each $i<r$ a local annihilation exponent $k_i$ and a
witness for $f_i^{k_i}m=0$.  Let $K=\max_i k_i$.  Raise each equation to the common exponent,
construct $d_i$ with $\sum d_i f_i^K=1$, and calculate
\[
                    m=\sum_i d_i f_i^K m=0.
\]
The induction invariant for the last fold records a witness that the accumulated partial sum equals
zero.

For existence, first homogenize the denominators of local sections to a common exponent $N$.
Take a maximum $K$ of the overlap-killing exponents and a safe $q=N+K+1$.  Form numerators $b_i$
so that the accepted overlap certificates imply
\[
                         f_j^q b_i=f_i^q b_j                \tag{14.4}
\]
for every $i,j$.  Construct $d_i$ with $\sum_i d_i f_i^q=1$ and output
\[
                              m=\sum_i d_i b_i.              \tag{14.5}\label{eq:glue-formula-object}
\]
For a fixed $j$, a bounded fold proves
\[
 f_j^q m=\sum_i d_i f_j^q b_i
          =\sum_i d_i f_i^q b_j=b_j.
\]
This equation is itself the localization equality witness between $m/1$ and the given local
section on $D(f_j)$.  A double-loop invariant covers all pairs $(i,j)$; no induction formula says
``for every localization''.

The two final statements should have types analogous to
\begin{lstlisting}[language={}] 
def principalCoverUniqueObjectProof :
  Prf PRAAx [] principalCoverUniqueFormula := ...

def principalCoverGlueObjectProof :
  Prf PRAAx [] principalCoverGlueFormula := ...
\end{lstlisting}

\subsection{Finite atlases and quasi-coherent data}

A finite affine atlas is a natural-number code containing:
\begin{itemize}
\item finitely presented chart rings $A_i$;
\item for each ordered pair $(i,j)$, finite principal-open refinements of the overlap;
\item forward and inverse localization homomorphisms with accepted composition witnesses;
\item triple-overlap cocycle witnesses;
\item a finite cover certificate wherever a family of principal opens is asserted to cover.
\end{itemize}
A quasi-coherent module contains chart modules $M_i$ and transition isomorphisms after the same
localizations.  Validity is one Boolean fold over the finite atlas.  The gluing theorem above gives
the equalizer law on every chart, and finite-family composition gives the corresponding law on the
atlas.  Morphisms, kernels, finite direct sums, and tensor products can then be added by certified
module operations.

This layer should be formalized as a library of certificate transformations, not as an internal
quantification over prime ideals or over all open subsets.  A separate semantic theorem in Lean may
interpret a valid atlas as a conventional glued scheme.

\subsection{The projective-line modules}

The $\mathbf P^1$ development uses sparse Laurent-polynomial codes.  Define
\[
       \delta_d(a,b)=a(t)-t^d b(t^{-1}).
\]
For a cocycle $c=\sum_e c_e t^e$, one bounded scan places:
\begin{align*}
 e\geq0 &\quad\hbox{in the first chart polynomial }a,\\
 e\leq d&\quad\hbox{in the second chart polynomial }b,\\
 d<e<0&\quad\hbox{in the obstruction }h.
\end{align*}
The invariant after processing a prefix of the finite support is
\[
       c_{<s}=\delta_d(a_s,b_s)+h_s,
       \qquad \operatorname{supp}(h_s)\subseteq\set{e:d<e<0}.
\]
This yields an object proof of \code{p1LaurentSplitFormula}.  Uniqueness is coefficientwise: neither
chart coboundary contains a monomial in the middle interval.  The $H^0$ theorem uses the same support
separation on a compatible pair.  In particular, at $d=0$ the global functions on $\mathbf P^1_R$
are exactly the constants in $R$.

\subsection{The projective-space homotopy}

Encode a subset $I\subseteq\{0,\ldots,n\}$ as a bit mask and an exponent vector as a signed list.
For one exponent vector $e$, let $N(e)$ be the negative-coordinate mask.  A component occurs at
$I$ exactly when $N(e)\subseteq I$.  The differential is a bounded alternating sum over deleted
bits.

When $N(e)$ is nonempty and not full, choose the least $i_*\notin N(e)$.  The homotopy inserts that
bit with sign determined by the number of smaller elements.  To prove
\[
                              dh+hd=\mathrm{id},             \tag{14.6}\label{eq:cech-homotopy-object}
\]
fix one output mask and scan all deletion positions.  The state records the unmatched term, the
paired contribution below $i_*$, and the paired contribution above $i_*$.  Terms not deleting
$i_*$ occur twice with opposite parity; the deletion of $i_*$ is the unique identity contribution.
The parity and bit-mask cases are decidable numerical predicates, so the induction is
quantifier-free.

The all-nonnegative case is the full simplex complex: the same insertion homotopy contracts the
positive degrees, while a coefficient projection identifies degree-zero constants.  The all-negative
case contains only the full mask.  A finite scan over the input cochain's support applies the
appropriate normalizer to each exponent vector.  This gives the object formula
\code{pnOdCechFormula} and the binomial-rank calculations by bounded counting of weak
compositions.

\subsection{Suggested module order}

A practical repository integration order is:
\begin{longtable}{@{}p{0.27\textwidth}p{0.65\textwidth}@{}}
\toprule
Module & Main contents and endpoint \\
\midrule
\endfirsthead
\toprule
Module & Main contents and endpoint \\
\midrule
\endhead
\code{AlgebraicGeometry/Codecs.lean} & Signed/list/polynomial/Laurent/matrix/subset codes and their PR graphs. \\
\code{AlgebraicGeometry/CertSetoid.lean} & \code{Valid}, \code{EqW}, witness transformers, bounded-fold congruence. \\
\code{AlgebraicGeometry/PresentedRing.lean} & Ideal-combination checker and quotient-operation witnesses. \\
\code{AlgebraicGeometry/Localization.lean} & Fraction codes and localization equivalence/congruence proofs. \\
\code{AlgebraicGeometry/PowerCover.lean} & Weak compositions, multinomial fold, and \code{powerCoverObjectProof}. \\
\code{AlgebraicGeometry/PrincipalDescent.lean} & Joint monicity and \code{principalCoverGlueObjectProof}. \\
\code{AlgebraicGeometry/Atlas.lean} & Finite affine atlases, refinements, morphisms, and module descent data. \\
\code{AlgebraicGeometry/ProjectiveLine.lean} & Two-chart atlas and Laurent normal forms for every $\cO(d)$. \\
\code{AlgebraicGeometry/ProjectiveCech.lean} & Exponent decomposition, insertion homotopy, and projective-space table. \\
\code{AlgebraicGeometry/NativeBridge.lean} & Applications of \code{pra_iff_native} to already constructed object derivations. \\
\bottomrule
\end{longtable}

\subsection{Completion criteria}

The formalization should be called complete only when all of the following hold:
\begin{enumerate}[label=\arabic*.]
\item the modules compile against the current \code{ConciseIPRA} sources;
\item the advertised formulas use the actual repository syntax and not a surrogate proof relation;
\item the principal theorems are inhabited \code{Prf PRAAx} values;
\item an axiom audit reports no \code{sorryAx}, custom axiom, or hidden native decision oracle;
\item finite executable tests agree with the object-level producers on a representative corpus;
\item the native translations are performed only after the preceding object proofs exist.
\end{enumerate}
The present archive supplies the mathematical derivations, numerical algorithms, host interfaces,
and this dependency order, but does not claim these completion criteria have already been met.
